%======================================================================

\subsection{Survey Reduction Mathematics}
\label{sec:survey_math}

This subsection details the mathematical transformations applied by the survey processing code to convert raw total station observations into geodetic coordinates. The procedure consists of three stages: (1)~conversion from polar observations to local ENU offsets, (2)~rotation from the local astronomic frame to ECEF, and (3)~conversion from ECEF back to geodetic coordinates.

\subsubsection{From Observations to Local ENU}

The total station measures the azimuth $A$ (from the calibrated north reference, Eq.~\ref{eq:azimuth_cal}), the elevation angle $e$ above the local horizontal, and the slant distance $S$ to each target. The corresponding offsets in the local East--North--Up frame are:
\begin{align}
\delta E &= S\,\cos e\,\sin A, \label{eq:obs2enu_E}\\
\delta N &= S\,\cos e\,\cos A, \label{eq:obs2enu_N}\\
\delta U &= S\,\sin e + h_{\text{inst}} - h_{\text{tgt}}, \label{eq:obs2enu_U}
\end{align}
where $h_{\text{inst}}$ is the instrument (trunnion axis) height above the base monument mark and $h_{\text{tgt}}$ is the reflector height above the target surface (zero for targets placed at the mirror surface; $\pm 0.18$\,m for the offset measurements described in Section~\ref{sec:total_station}). The slant distance $S$ has already been corrected for the 30\,mm reflector offset (for standard prism reflectors) and converted from feet to meters.

These offsets are relative to the base monument and are expressed in the \emph{astronomic} (gravity-aligned) local frame, because the total station's compensator aligns to the plumb line.

\subsubsection{From Local ENU to ECEF (with Deflection of the Vertical)}

Converting the local ENU offsets to the geocentric ECEF frame proceeds in two steps.

\paragraph{Step 1: Astronomic to geodetic ENU.} The deflection-of-vertical rotation (Eq.~\ref{eq:deflection_rotation}) is applied to convert from the astronomic ENU frame (in which the instrument operates) to the geodetic ENU frame (aligned with the ellipsoid normal):
\begin{equation}
\begin{pmatrix} \delta E_g \\ \delta N_g \\ \delta U_g \end{pmatrix}
=
\begin{pmatrix}
1 & 0 & \eta \\
0 & 1 & -\xi \\
-\eta & \xi & 1
\end{pmatrix}
\begin{pmatrix} \delta E \\ \delta N \\ \delta U \end{pmatrix},
\label{eq:astro_to_geod}
\end{equation}
where $\xi$ and $\eta$ are the deflection parameters from Table~\ref{tab:deflection}, converted to radians.

\paragraph{Step 2: Geodetic ENU to ECEF.} The standard rotation from the geodetic local ENU frame to ECEF uses the geodetic latitude $\varphi$ and longitude $\lambda$ of the base monument:
\begin{align}
\Delta X &= -\sin\lambda\;\delta E_g - \sin\varphi\cos\lambda\;\delta N_g + \cos\varphi\cos\lambda\;\delta U_g, \label{eq:enu2ecef_X}\\
\Delta Y &= +\cos\lambda\;\delta E_g - \sin\varphi\sin\lambda\;\delta N_g + \cos\varphi\sin\lambda\;\delta U_g, \label{eq:enu2ecef_Y}\\
\Delta Z &= \phantom{+\cos\lambda\;\delta E_g} + \cos\varphi\;\delta N_g + \sin\varphi\;\delta U_g. \label{eq:enu2ecef_Z}
\end{align}
In matrix form, defining $\mathbf{R}_{\text{ENU}}$ as the rotation matrix with rows given by the coefficients above:
\begin{equation}
\begin{pmatrix} \Delta X \\ \Delta Y \\ \Delta Z \end{pmatrix}
=
\begin{pmatrix}
-\sin\lambda & -\sin\varphi\cos\lambda & \cos\varphi\cos\lambda \\
+\cos\lambda & -\sin\varphi\sin\lambda & \cos\varphi\sin\lambda \\
0 & +\cos\varphi & +\sin\varphi
\end{pmatrix}
\begin{pmatrix} \delta E_g \\ \delta N_g \\ \delta U_g \end{pmatrix}.
\label{eq:enu2ecef_matrix}
\end{equation}

The target's ECEF coordinates are then:
\begin{equation}
\begin{pmatrix} X \\ Y \\ Z \end{pmatrix}_{\!\text{target}}
=
\begin{pmatrix} X \\ Y \\ Z \end{pmatrix}_{\!\text{base}}
+
\begin{pmatrix} \Delta X \\ \Delta Y \\ \Delta Z \end{pmatrix}.
\label{eq:target_ecef}
\end{equation}
The base monument ECEF coordinates $(X, Y, Z)_{\text{base}}$ are computed from the geodetic monument coordinates using Eqs.~\eqref{eq:geod2xyz_X}--\eqref{eq:geod2xyz_Z}.

\subsubsection{From ECEF to Geodetic}

The target ECEF coordinates $(X, Y, Z)_{\text{target}}$ are converted back to geodetic coordinates $(\varphi, \lambda, h)_{\text{target}}$ using Bowring's iterative method (Eqs.~\ref{eq:xyz2geod_lam}--\ref{eq:bowring_h}). For the short baselines in this survey ($< 1$\,km except CLF--MD at $\sim\!20$\,km), a single Bowring iteration achieves sub-millimeter convergence.

The final geodetic coordinates inherit the reference frame of the input monument coordinates: if the base monument is in ITRF00, then the computed target position is in ITRF00. This is the mechanism by which the 0.73\,m NAD\,83/ITRF height discrepancy propagates from the monuments to the telescope mirrors, as discussed throughout this report. The relative geometry (mirror-to-mirror distances) is independent of the reference frame because the ENU-to-ECEF-to-geodetic round-trip is a rigid transformation.


%======================================================================
% BIBLIOGRAPHY
%======================================================================

\begin{thebibliography}{16}

\bibitem{misra2006}
Misra, P.\ \& Enge, P.,
\textit{Global Positioning System: Signals, Measurements, and Performance},
2nd ed., Ganga-Jamuna Press, 2006.

\bibitem{hofmann2008}
Hofmann-Wellenhof, B., Lichtenegger, H., \& Wasle, E.,
\textit{GNSS: Global Navigation Satellite Systems},
Springer, 2008.

\bibitem{klobuchar1987}
Klobuchar, J.\,A.,
``Ionospheric Time-Delay Algorithm for Single-Frequency GPS Users,''
\textit{IEEE Trans.\ Aerospace and Electronic Systems}, AES-23(3), 325--331, 1987.

\bibitem{coster2008}
Coster, A.\ \& Komjathy, A.,
``Space Weather and the Global Positioning System,''
\textit{Space Weather}, 6, S06D04, 2008.

\bibitem{ovstedal2002}
{\O}vstedal, O.,
``Absolute Positioning with Single-Frequency GPS Receivers,''
\textit{GPS Solutions}, 6(3), 151--160, 2002.

\bibitem{mannucci1998}
Mannucci, A.\,J., et al.,
``A global mapping technique for GPS-derived ionospheric TEC measurements,''
\textit{Radio Science}, 33(3), 565--582, 1998.

\bibitem{soler2004}
Soler, T.\ \& Snay, R.\,A.,
``Transforming Positions and Velocities between ITRF2000 and NAD\,83,''
\textit{J.\ Surveying Engineering}, 130(2), 49--55, 2004.

\bibitem{soler2003}
Soler, T., et al.,
``OPUS: Online Positioning User Service,''
\textit{Professional Surveyor}, 23(5), 2003.

\bibitem{roman2012}
Roman, D.\,R.,
``A National Geoid Model for NAD\,83,''
NOAA/NGS, 2012.

\bibitem{nga2014}
National Geospatial-Intelligence Agency,
``DoD World Geodetic System 1984,''
NGA.STND.0036, 4th ed., 2014.

\bibitem{bowring1976}
Bowring, B.\,R.,
``Transformation from spatial to geographical coordinates,''
\textit{Survey Review}, 23(181), 323--327, 1976.

\bibitem{torge2012}
Torge, W.\ \& M{\"u}ller, J.,
\textit{Geodesy}, 4th ed., de Gruyter, 2012.

\bibitem{gesch2002}
Gesch, D., et al.,
``The National Elevation Dataset,''
\textit{Photogrammetric Engineering and Remote Sensing}, 68(1), 5--11, 2002.

\bibitem{tokuno}
Tokuno, H., et al.,
unpublished (cited in Nakazawa 2020).

\bibitem{matsuzawa2025}
Matsuzawa, A.,
``MD Opt-copter Results,''
TA-ALL Collaboration Meeting, June 2025.

\bibitem{nakazawa2020}
Nakazawa, A.,
``BRM Opt-copter Results,''
TA Internal Report, 2020.

\end{thebibliography}
